% Options for packages loaded elsewhere
\PassOptionsToPackage{unicode}{hyperref}
\PassOptionsToPackage{hyphens}{url}
%
\documentclass[
]{article}
\usepackage{amsmath,amssymb}
\usepackage{iftex}
\ifPDFTeX
  \usepackage[T1]{fontenc}
  \usepackage[utf8]{inputenc}
  \usepackage{textcomp} % provide euro and other symbols
\else % if luatex or xetex
  \usepackage{unicode-math} % this also loads fontspec
  \defaultfontfeatures{Scale=MatchLowercase}
  \defaultfontfeatures[\rmfamily]{Ligatures=TeX,Scale=1}
\fi
\usepackage{lmodern}
\ifPDFTeX\else
  % xetex/luatex font selection
\fi
% Use upquote if available, for straight quotes in verbatim environments
\IfFileExists{upquote.sty}{\usepackage{upquote}}{}
\IfFileExists{microtype.sty}{% use microtype if available
  \usepackage[]{microtype}
  \UseMicrotypeSet[protrusion]{basicmath} % disable protrusion for tt fonts
}{}
\makeatletter
\@ifundefined{KOMAClassName}{% if non-KOMA class
  \IfFileExists{parskip.sty}{%
    \usepackage{parskip}
  }{% else
    \setlength{\parindent}{0pt}
    \setlength{\parskip}{6pt plus 2pt minus 1pt}}
}{% if KOMA class
  \KOMAoptions{parskip=half}}
\makeatother
\usepackage{xcolor}
\usepackage[margin=1in]{geometry}
\usepackage{longtable,booktabs,array}
\usepackage{calc} % for calculating minipage widths
% Correct order of tables after \paragraph or \subparagraph
\usepackage{etoolbox}
\makeatletter
\patchcmd\longtable{\par}{\if@noskipsec\mbox{}\fi\par}{}{}
\makeatother
% Allow footnotes in longtable head/foot
\IfFileExists{footnotehyper.sty}{\usepackage{footnotehyper}}{\usepackage{footnote}}
\makesavenoteenv{longtable}
\usepackage{graphicx}
\makeatletter
\def\maxwidth{\ifdim\Gin@nat@width>\linewidth\linewidth\else\Gin@nat@width\fi}
\def\maxheight{\ifdim\Gin@nat@height>\textheight\textheight\else\Gin@nat@height\fi}
\makeatother
% Scale images if necessary, so that they will not overflow the page
% margins by default, and it is still possible to overwrite the defaults
% using explicit options in \includegraphics[width, height, ...]{}
\setkeys{Gin}{width=\maxwidth,height=\maxheight,keepaspectratio}
% Set default figure placement to htbp
\makeatletter
\def\fps@figure{htbp}
\makeatother
\setlength{\emergencystretch}{3em} % prevent overfull lines
\providecommand{\tightlist}{%
  \setlength{\itemsep}{0pt}\setlength{\parskip}{0pt}}
\setcounter{secnumdepth}{-\maxdimen} % remove section numbering
\ifLuaTeX
  \usepackage{selnolig}  % disable illegal ligatures
\fi
\IfFileExists{bookmark.sty}{\usepackage{bookmark}}{\usepackage{hyperref}}
\IfFileExists{xurl.sty}{\usepackage{xurl}}{} % add URL line breaks if available
\urlstyle{same}
\hypersetup{
  pdftitle={321 Final Project},
  pdfauthor={Ariel Lin and Naomi Le Mouel},
  hidelinks,
  pdfcreator={LaTeX via pandoc}}

\title{321 Final Project}
\author{Ariel Lin and Naomi Le Mouel}
\date{2024-03-13}

\begin{document}
\maketitle

\hypertarget{introduction}{%
\section{Introduction}\label{introduction}}

Social media has a powerful presence in American lives. According to
data from Statista, there are over 300 million social media users in the
United States. Media platforms are used for many purposes, including
entertainment, socialization, business, education, and news. Around half
of the US population gets their news information from social media,
according to a 2021 study conducted by the Pew Research Center. With so
many media influencers and users, there is a lot of room for
misinformation, fake news, propaganda, and online discourse, which can
often concern political topics such as presidential elections.

In this report, we examine the potential relationship between Americans'
satisfaction with U.S. democracy and their opinions on social media's
impact on democracy. The data used in this report is from the Pew
Research Center's American Trends Panel (ATP) Wave 105, which surveyed
3,581 U.S. adults from March 21 to 27, 2022, through random sampling of
residential addresses and weighted to represent the U.S. population by
gender, race, and other demographics.

\hypertarget{theory-and-hypotheses}{%
\section{Theory and Hypotheses}\label{theory-and-hypotheses}}

Using the available ATP Wave 105 data, we want to answer this question:
what effect has social media had on democracy in the eyes of adults in
the United States? In other words, how is American satisfaction with
U.S. democracy related to their opinions on social media's impact on
democracy? Our report draws from a broader Pew Research Center report by
Wike et al.~published in December 2022 that observes the trend in
attitudes toward social media and its effects on democracy in several
countries across the globe. The study used the ATP Wave 105 data for the
U.S. and randomized phone surveys for non-U.S. countries and found that
social media was perceived as ``mostly good'' by most countries in
general, except the United States. Attitudes toward social media's
effect on democracy are further explored more specifically through
survey questions asking about social media's impact on positive
outcomes, such as better informing people, and negative outcomes, such
as manipulating people.

Here, we focus on the United States and make the causal assumption that
low American satisfaction with democracy in the US should have a
positive correlation with the belief that social media harms democracy
(and vice versa). We reason that the people who believe social media is
``bad'' for democracy are also unsatisfied with the US democracy because
social media contributed to the unsatisfactory democracy.

Based on our assumption, two hypotheses are considered:

\begin{itemize}
\item
  \(H_0\) (null hypothesis) - American adult's belief that social media
  is ``bad'' or ``good'' for democracy does not affect
  satisfaction/dissatisfaction with democracy.
\item
  \(H_1\) (alternative hypothesis) - American adult's belief that social
  media is ``bad'' or ``good'' for democracy does affect
  satisfaction/dissatisfaction with democracy.
\end{itemize}

We would also like to consider a few possible heterogeneous effects:

\begin{itemize}
\item
  Posting social/political issues on social media: Do people who post
  social and political issues on social media believe social media has a
  ``good'' impact on democracy? (Believing they can use the platform to
  educate others on issues etc.)
\item
  Ideology: are liberal respondents more likely to be dissatisfied with
  the U.S. democracy? Are moderates more likely to be satisfied?
\item
  Age: do older people perceive social media's impacts on democracy more
  negatively than younger people?
\item
  Race: Are non-White minorities less satisfied with the U.S. democracy?
\end{itemize}

\hypertarget{data-and-visual-analysis}{%
\section{Data and Visual Analysis}\label{data-and-visual-analysis}}

\hypertarget{data}{%
\subsection{Data}\label{data}}

The ATP Wave 105 dataset conducted by the Pew Research Center contains
3,581 observations representing survey questionnaire responses of
individual U.S. adults and 159 variables corresponding to each question
on the survey. Our study is observational since we are looking at
relationships between responses to relevant questions on the ATP Wave
105 questionnaire.

We selected the following questions and variables to focus on for our
project:

\texttt{SATISFIED\_DEMOCRACY\_w105}: How satisfied are you with the way
democracy is working in the U.S?

\begin{enumerate}
\def\labelenumi{\arabic{enumi})}
\item
  Very satisfied
\item
  Somewhat satisfied
\item
  Not too satisfied
\item
  Not at all satisfied
\end{enumerate}

\texttt{F\_PARTY\_FINAL}: As of today do you lean more to

\begin{enumerate}
\def\labelenumi{\arabic{enumi})}
\item
  Republican
\item
  Democrat
\item
  Independent
\item
  Something else
\end{enumerate}

\texttt{SNSUSE}: Do you use social media sites like Facebook, Twitter,
or Instagram?

\begin{enumerate}
\def\labelenumi{\arabic{enumi})}
\item
  Yes, I use social media sites
\item
  No, I do not use social media sites
\end{enumerate}

\texttt{SOCIAL\_DEMOCRACY\_w150}: Overall, when you consider all the
advantages and disadvantages of social media, would you say social media
has been

\begin{enumerate}
\def\labelenumi{\arabic{enumi})}
\item
  More of a GOOD THING for democracy in the U.S.
\item
  More of a BAD THING for democracy in the U.S.
\end{enumerate}

\texttt{SMIMPACT\_BATT}: In general, do you think access to the internet
and social media has made people more or less

\texttt{SMIMPACT\_INFORMED\_w105}: Informed about current events in our
country?

\texttt{SMIMPACT\_DIVIDED\_w105}: Divided in their political opinions?

\texttt{SMIMPACT\_MANIPULATE\_w105}: Easy to manipulate with false
information and rumors?

\texttt{SMIMPACT\_DISCOURSE\_w105}: Civil in the way they talk about
politics?

\begin{enumerate}
\def\labelenumi{\arabic{enumi})}
\item
  More
\item
  Less
\item
  Not had much impact
\end{enumerate}

We also included variables from the dataset that were not associated
with any particular question in the survey:

\texttt{F\_AGECAT}: age category of the respondents

\begin{enumerate}
\def\labelenumi{\arabic{enumi})}
\item
  65+
\item
  50-64
\item
  30-49
\item
  18-29
\end{enumerate}

\texttt{F\_IDEO}: ideology of the respondents

\begin{enumerate}
\def\labelenumi{\arabic{enumi})}
\item
  Very conservative
\item
  Conservative
\item
  Liberal
\item
  Moderate
\item
  Very liberal
\end{enumerate}

\texttt{F\_RACECMB}: race the respondent identifies with

\begin{enumerate}
\def\labelenumi{\arabic{enumi})}
\item
  White
\item
  Asian or Asian-American
\item
  Black or African-American
\item
  Mixed Race
\item
  Or some other race
\end{enumerate}

The original response data ATP Wave 105 dataset was qualitative, i.e.,
stored as words and phrases. To better observe relationships between the
variables, we created new variables that quantified the data by
assigning numerical values representing the responses, which we stored
into a new, smaller dataset of 11 quantitative variables instead of 159.
We also created a dataset with 22 variables that included the original
qualitative responses for the variables. A summary of the variables of
our quantitative modified dataset is shown below.

\begin{longtable}[]{@{}
  >{\raggedright\arraybackslash}p{(\columnwidth - 2\tabcolsep) * \real{0.3118}}
  >{\raggedright\arraybackslash}p{(\columnwidth - 2\tabcolsep) * \real{0.6882}}@{}}
\toprule\noalign{}
\begin{minipage}[b]{\linewidth}\raggedright
Name
\end{minipage} & \begin{minipage}[b]{\linewidth}\raggedright
Description
\end{minipage} \\
\midrule\noalign{}
\endhead
\bottomrule\noalign{}
\endlastfoot
\texttt{satisfied} & Satisfaction of respondents with Democracy in the
US on a scale of 1 (not at all satisfied) to 4 (very satisfied) \\
\texttt{socialuse} & Whether or not respondents use social media sites
(yes, 1 or no, 0) \\
\texttt{social\_demo} & Whether respondents think social media is more
of a GOOD (1) thing or a BAD (0) thing for democracy \\
\texttt{posting} & Whether or not respondents posted about political or
social issues on social media on a scale of 1(often) to 4(never) \\
\texttt{smp\_informed} & If respondents think access to the internet and
social media makes people more or less informed about current events in
our country (more, 1, less, 2, not much impact, 3) \\
\texttt{smn\_manip} & If respondents think access to the internet and
social media makes people more or less manipulated with false
information and rumors (more, 1, less, 2, not much impact, 3) \\
\texttt{smn\_divided} & If respondents think access to the internet and
social media makes people more or less divided in their political
opinions (more, 1, less, 2, not much impact, 3) If respondents think
access to the internet and social media makes people more or less civil
in the way they talk about politics (more, 1, less, 2, not much impact,
3) \\
\texttt{ideology} & The ideology of the respondents on a scale of 1
(liberal) to 3 (conservative) \\
\texttt{agecat} & The age category the respondent falls under on a scale
of 1 to 4 (18-29, 30-49, 50-64, 65+) \\
\texttt{race} & The race of the respondent (mixed race 1, Asian or Asian
american 2, some other race 3, Black or African American 4, white 5) \\
\end{longtable}

We omitted missing values when conducting our the majority of our data
analyses, since most variables have only small proportion of missing
values (all under 5 percent). There are 2426 observations that do not
contain missing values. The \texttt{posting} variable is the only
exception, which we assume may be due to some respondents being
uncomfortable disclosing whether or not they post political and social
issues on social media. (\emph{See Table 1})

\begin{table}[H] \centering 
  \caption{Proportion of Missing Values} 
  \label{} 
\begin{tabular}{@{\extracolsep{5pt}} ccc} 
\\[-1.8ex]\hline 
\hline \\[-1.8ex] 
Variable & Num\_NAs & Prop\_NAs \\ 
\hline \\[-1.8ex] 
posting & $1,008$ & $0.281$ \\ 
social\_demo & $64$ & $0.018$ \\ 
satisfied & $20$ & $0.006$ \\ 
ideology & $72$ & $0.020$ \\ 
socialuse & $144$ & $0.040$ \\ 
agecat & $13$ & $0.004$ \\ 
race & $63$ & $0.018$ \\ 
smn\_manip & $23$ & $0.006$ \\ 
smn\_divided & $38$ & $0.011$ \\ 
smp\_informed & $25$ & $0.007$ \\ 
smp\_civil & $31$ & $0.009$ \\ 
\hline \\[-1.8ex] 
\end{tabular} 
\end{table}

The mean of 0.33 for the \texttt{social\_demo} variable suggests that a
third of U.S. adults think social media is ``good'' for democracy and
the majority think social media is ``bad'' for democracy. Most U.S.
adults are also not very satisfied with the U.S. democracy, based on the
low mean of 2.23 for \texttt{satisfied}. The averages of the specific
effects (\texttt{smp\_inform}, \texttt{smp\_divided},
\texttt{smn\_manip}) of social media on democracy are fairly similar,
implying that for both positive and negative effects, U.S. adults
generally think social media has a ``more of an impact'' on those
effects (manipulation, political division, information) and ``less of an
impact'' on civil discussions about politics (\texttt{smp\_civil}). The
majority of respondents in the dataset are between 30 and 64 years of
age, white, use social media, have a moderate ideology, and rarely post
about social and political issues on social media. (\emph{See Table 2})

\begin{table}[H] \centering 
  \caption{Descriptive statistics} 
  \label{} 
\begin{tabular}{@{\extracolsep{5pt}}lccccc} 
\\[-1.8ex]\hline 
\hline \\[-1.8ex] 
Statistic & \multicolumn{1}{c}{Mean} & \multicolumn{1}{c}{St. Dev.} & \multicolumn{1}{c}{Min} & \multicolumn{1}{c}{Median} & \multicolumn{1}{c}{Max} \\ 
\hline \\[-1.8ex] 
posting & 3.05 & 0.91 & 1 & 3 & 4 \\ 
social\_demo & 0.33 & 0.47 & 0 & 0 & 1 \\ 
satisfied & 2.23 & 0.80 & 1 & 2 & 4 \\ 
ideology & 2.08 & 0.78 & 1 & 2 & 3 \\ 
socialuse & 0.75 & 0.43 & 0 & 1 & 1 \\ 
agecat & 2.67 & 0.99 & 1 & 3 & 4 \\ 
race & 4.51 & 1.03 & 1 & 5 & 5 \\ 
smn\_manip & 1.18 & 0.52 & 1 & 1 & 3 \\ 
smn\_divided & 1.30 & 0.68 & 1 & 1 & 3 \\ 
smp\_informed & 1.50 & 0.74 & 1 & 1 & 3 \\ 
smp\_civil & 2.00 & 0.52 & 1 & 2 & 3 \\ 
\hline \\[-1.8ex] 
\end{tabular} 
\end{table}

Based on our correlation matrix, the variable \texttt{social\_demo}
exhibits a positive correlation of 0.19 with the variable
\texttt{satisfied}. Additionally, \texttt{posting} demonstrates a small
positive correlation of 0.06 with satisfaction with the U.S. democracy,
while \texttt{ideology} displays a negative correlation of -0.20,
suggesting that more respondents with conservative ideologies are more
likely to be less satisfied with the U.S. democracy. (\emph{See Table
3})

\begin{table}[H] \centering 
  \caption{Correlation matrix} 
  \label{} 
\small 
\begin{tabular}{@{\extracolsep{-10pt}} cccccccccccc} 
\\[-1.8ex]\hline 
\hline \\[-1.8ex] 
 & posting & social\_demo & satisfied & ideology & socialuse & agecat & race & smn\_manip & smn\_divided & smp\_informed & smp\_civil \\ 
\hline \\[-1.8ex] 
posting & $1$ & $$-$0.14$ & $0.06$ & $0.08$ & $$ & $$-$0.05$ & $0.03$ & $$-$0.02$ & $0.03$ & $0.07$ & $0.10$ \\ 
social\_demo & $$-$0.14$ & $1$ & $0.19$ & $$-$0.16$ & $$ & $$-$0.11$ & $$-$0.14$ & $0.12$ & $0.16$ & $$-$0.21$ & $$-$0.04$ \\ 
satisfied & $0.06$ & $0.19$ & $1$ & $$-$0.20$ & $$ & $$-$0.01$ & $$-$0.07$ & $$-$0.01$ & $0.03$ & $$-$0.11$ & $$-$0.06$ \\ 
ideology & $0.08$ & $$-$0.16$ & $$-$0.20$ & $1$ & $$ & $0.22$ & $0.11$ & $0.01$ & $$-$0.01$ & $0.07$ & $0$ \\ 
socialuse & $$ & $$ & $$ & $$ & $1$ & $$ & $$ & $$ & $$ & $$ & $$ \\ 
agecat & $$-$0.05$ & $$-$0.11$ & $$-$0.01$ & $0.22$ & $$ & $1$ & $0.17$ & $$-$0.02$ & $$-$0.03$ & $0.13$ & $0.01$ \\ 
race & $0.03$ & $$-$0.14$ & $$-$0.07$ & $0.11$ & $$ & $0.17$ & $1$ & $$-$0.04$ & $$-$0.09$ & $0.02$ & $0.04$ \\ 
smn\_manip & $$-$0.02$ & $0.12$ & $$-$0.01$ & $0.01$ & $$ & $$-$0.02$ & $$-$0.04$ & $1$ & $0.35$ & $0.14$ & $0.14$ \\ 
smn\_divided & $0.03$ & $0.16$ & $0.03$ & $$-$0.01$ & $$ & $$-$0.03$ & $$-$0.09$ & $0.35$ & $1$ & $0.08$ & $0.19$ \\ 
smp\_informed & $0.07$ & $$-$0.21$ & $$-$0.11$ & $0.07$ & $$ & $0.13$ & $0.02$ & $0.14$ & $0.08$ & $1$ & $0.15$ \\ 
smp\_civil & $0.10$ & $$-$0.04$ & $$-$0.06$ & $0$ & $$ & $0.01$ & $0.04$ & $0.14$ & $0.19$ & $0.15$ & $1$ \\ 
\hline \\[-1.8ex] 
\end{tabular} 
\end{table}

\hypertarget{exploratory-data-analysis}{%
\subsection{Exploratory Data Analysis}\label{exploratory-data-analysis}}

We begin by observing the distribution of U.S. adults' satisfaction with
U.S. democracy.

For the entire dataset, including the small number of respondents who do
not use social media, the percentage for each response of the total
number of respondents is displayed in Table 4. The vast majority of the
responses are around ``2'' and ``3'', which reflects that most
respondents are either ``Not too satisfied'' (41.7\%) or only ``Somewhat
satisfied'' (34.5\%) with the U.S. democracy. Note: the last row
represents the percentage of missing values, which are under 1 percent
for both tables below.

\begin{table}[H] \centering 
  \caption{Distribution of Satisfaction with U.S. Democracy} 
  \label{} 
\begin{tabular}{@{\extracolsep{5pt}} cc} 
\\[-1.8ex]\hline 
\hline \\[-1.8ex] 
Satisfaction & \% of Respondents \\ 
\hline \\[-1.8ex] 
$1$ & $19.44$ \\ 
$2$ & $41.66$ \\ 
$3$ & $34.49$ \\ 
$4$ & $3.85$ \\ 
$$ & $0.56$ \\ 
\hline \\[-1.8ex] 
\end{tabular} 
\end{table}

Subsetting the data for only those respondents who \emph{do} use social
media, we still see a similar trend, with a slightly higher percentage
for ``Not too satisfied'' (about 2 percentage points higher) with U.S.
democracy.

\begin{table}[H] \centering 
  \caption{Distribution of Satisfaction with U.S. Democracy among Social Media Users} 
  \label{} 
\begin{tabular}{@{\extracolsep{5pt}} cc} 
\\[-1.8ex]\hline 
\hline \\[-1.8ex] 
Satisfaction & \% of Respondents \\ 
\hline \\[-1.8ex] 
$1$ & $18.99$ \\ 
$2$ & $43.68$ \\ 
$3$ & $33.19$ \\ 
$4$ & $3.68$ \\ 
$$ & $0.46$ \\ 
\hline \\[-1.8ex] 
\end{tabular} 
\end{table}

Plotting the percentages by category of satisfaction for all respondents
and for only respondents that use social media, we again see that the
distributions are very similar, with the majority of responses in the
``Not too satisfied'' with U.S. democracy category.

\includegraphics[width=0.5\linewidth]{321FINALv2_files/figure-latex/unnamed-chunk-7-1}
\includegraphics[width=0.5\linewidth]{321FINALv2_files/figure-latex/unnamed-chunk-7-2}

Now we compare the satisfaction with U.S. democracy of respondents who
believe social media is more of a bad thing for democracy with that of
respondents who believe social media is more of a good thing. We subset
the data into these two groups and observe the trends in the percentage
of responses for each satisfaction category.

Based on Table 7, the majority of responses (43.61\%) among U.S. adults
that believe social media is more of a \emph{bad} thing for democracy
are ``2'', meaning they are ``Not too satisfied'' with U.S. democracy.

\begin{table}[H] \centering 
  \caption{Satisfaction with U.S. Democracy among Respondents Believing Social Media is Bad for Democracy} 
  \label{} 
\begin{tabular}{@{\extracolsep{5pt}} cc} 
\\[-1.8ex]\hline 
\hline \\[-1.8ex] 
Satisfaction & \% of Respondents \\ 
\hline \\[-1.8ex] 
$1$ & $23.67$ \\ 
$2$ & $43.61$ \\ 
$4$ & $2.56$ \\ 
$3$ & $29.77$ \\ 
$$ & $0.38$ \\ 
\hline \\[-1.8ex] 
\end{tabular} 
\end{table}

In contrast, the majority of responses (43.71\%) among U.S. adults that
believe social media is more of a \emph{good} thing for democracy are
``3'', meaning they are ``somewhat satisfied'' with U.S. democracy.

\begin{table}[H] \centering 
  \caption{Satisfaction with U.S. Democracy among Respondents Believing Social Media is Good for Democracy} 
  \label{} 
\begin{tabular}{@{\extracolsep{5pt}} cc} 
\\[-1.8ex]\hline 
\hline \\[-1.8ex] 
Satisfaction & \% of Respondents \\ 
\hline \\[-1.8ex] 
$2$ & $38.27$ \\ 
$3$ & $43.71$ \\ 
$4$ & $6.38$ \\ 
$1$ & $11.31$ \\ 
$$ & $0.34$ \\ 
\hline \\[-1.8ex] 
\end{tabular} 
\end{table}

Comparing the percentage of respondents who responded ``Not too
satisfied'' for the two tables, the percentage of respondents who
believe social media is more of a good thing for democracy is 5.35
percentage points lower than that of respondents who believe social
media is more of a bad thing for democracy.

The bar plots help illustrate our observations. Comparing the two bar
plots, we see that the distribution for respondents believing social
media is bad for democracy is slightly more skewed right, while the
distribution for respondents believing social media is good for
democracy is slightly more skewed to the left. The skewed distributions
indicate that those who believe social media is good for democracy are
more likely to be satisfied with U.S. democracy and those who believe
social media is bad for democracy are less likely satisfied with the
U.S. democracy.

\includegraphics[width=0.5\linewidth]{321FINALv2_files/figure-latex/unnamed-chunk-11-1}
\includegraphics[width=0.5\linewidth]{321FINALv2_files/figure-latex/unnamed-chunk-11-2}

For more specific impacts of social media, the Pew Research study
considers the variables \texttt{smp\_informed} and \texttt{smp\_civil}
good impacts of social media and \texttt{smn\_manip} and
\texttt{smn\_divided} bad impacts of social media on democracy (see the
table on page 4 for descriptions of the variables). We presume that
respondents who believe social media is good for democracy will more
likely believe that social media has ``more of an impact'' (coded as a 1
in our quantitative variable) on the ``good'' variables and less of an
impact or no impact (2 or 3) on the ``bad'' variables. The opposite
might be true for respondents believing social media has an overall bad
impact on democracy.

\begin{table}[H] \centering 
  \caption{Social Media's Specific Impacts for Respondents Believing Social Media is Good for Democracy} 
  \label{} 
\begin{tabular}{@{\extracolsep{5pt}} cc} 
\\[-1.8ex]\hline 
\hline \\[-1.8ex] 
Impact & Mean \\ 
\hline \\[-1.8ex] 
smp\_civil & $1.97$ \\ 
smp\_informed & $1.28$ \\ 
smn\_manip & $1.27$ \\ 
smn\_divided & $1.45$ \\ 
\hline \\[-1.8ex] 
\end{tabular} 
\end{table}

\begin{table}[H] \centering 
  \caption{Social Media's Specific Impacts for Respondents Believing Social Media is Bad for Democracy} 
  \label{} 
\begin{tabular}{@{\extracolsep{5pt}} cc} 
\\[-1.8ex]\hline 
\hline \\[-1.8ex] 
Impact & Mean \\ 
\hline \\[-1.8ex] 
smp\_civil & $2.02$ \\ 
smp\_informed & $1.62$ \\ 
smn\_manip & $1.13$ \\ 
smn\_divided & $1.22$ \\ 
\hline \\[-1.8ex] 
\end{tabular} 
\end{table}

According to the values in the table, respondents believing social media
is good for democracy think social media's ability to manipulate is the
strongest impact and much less of an impact on informing people on
current events going on in the country. We do, however, see a slight
difference for respondents who believe social media is bad for democracy
- the means for the ``good'' impacts (\texttt{smp\_civil} and
\texttt{smp\_informed}) are visibly greater than the means for the
``bad'' impacts, supporting the idea that those who believe social media
is bad for democracy think social media's ``bad'' abilities to
manipulate and divide is stronger than its ``good'' abilities to inform
and promote civil discourse. In general, the means for all four
variables are somewhat similar, suggesting that respondents' belief in
whether social media is overall good or bad for democracy is not largely
based on the defined ``good'' and ``bad'' specific impacts. (\emph{See
Tables 9 and 10})

We are also interested in observing possible differences in satisfaction
and belief that social media is good or bad for democracy by race,
political ideology, age category, and whether the respondents post
social and political issues on social media.

The mean satisfaction by race for respondents believing social media is
good for democracy are overall similar (mostly not too satisfied or only
somewhat satisfied with U.S. democracy) and slightly higher than the
satisfaction of respondents believing social media is bad. Those who are
mixed race and believe social media is bad for democracy have the lowest
satisfaction (mean of 1.98) with U.S. democracy.

\begin{table}[ht]
\centering
\caption{Respondents Believing Social Media is Good
             and Mean Satisfaction by Race} 
\begin{tabular}{lr}
  \hline
Race & Mean Satisfaction \\ 
  \hline
Asian or Asian-American & 2.51 \\ 
  Black or African-American & 2.40 \\ 
  Mixed Race & 2.45 \\ 
  Or some other race & 2.55 \\ 
  White & 2.38 \\ 
   \hline
\end{tabular}
\end{table}
\begin{table}[ht]
\centering
\caption{Respondents Believing Social Media is Bad
             and Mean Satisfaction by Race} 
\begin{tabular}{lr}
  \hline
Race & Mean Satisfaction \\ 
  \hline
Asian or Asian-American & 2.25 \\ 
  Black or African-American & 2.31 \\ 
  Mixed Race & 1.98 \\ 
  Or some other race & 2.30 \\ 
  White & 2.07 \\ 
   \hline
\end{tabular}
\end{table}

\begin{table}[ht]
\centering
\caption{Respondents Believing Social Media is Good
             and Mean Satisfaction by Ideology} 
\begin{tabular}{rr}
  \hline
Ideology & Mean Satisfaction \\ 
  \hline
1.00 & 2.39 \\ 
  2.00 & 2.57 \\ 
  3.00 & 2.16 \\ 
   \hline
\end{tabular}
\end{table}
\begin{table}[ht]
\centering
\caption{Respondents Believing Social Media is Bad
             and Mean Satisfaction by Ideology} 
\begin{tabular}{rr}
  \hline
Ideology & Mean Satisfaction \\ 
  \hline
1.00 & 2.30 \\ 
  2.00 & 2.22 \\ 
  3.00 & 1.86 \\ 
   \hline
\end{tabular}
\end{table}
\begin{table}[ht]
\centering
\caption{Respondents Believing Social Media is Good
             and Mean Satisfaction by Age Category} 
\begin{tabular}{lr}
  \hline
Age Category & Mean Satisfaction \\ 
  \hline
18-29 & 2.27 \\ 
  30-49 & 2.41 \\ 
  50-64 & 2.54 \\ 
  65+ & 2.37 \\ 
   \hline
\end{tabular}
\end{table}
\begin{table}[ht]
\centering
\caption{Respondents Believing Social Media is Bad
             and Mean Satisfaction by Age Category} 
\begin{tabular}{lr}
  \hline
Age Category & Mean Satisfaction \\ 
  \hline
18-29 & 2.09 \\ 
  30-49 & 2.14 \\ 
  50-64 & 2.07 \\ 
  65+ & 2.07 \\ 
   \hline
\end{tabular}
\end{table}

\begin{table}[ht]
\centering
\caption{Respondents Believing Social Media is Good
             and Mean Satisfaction by How Often They Post Issues} 
\begin{tabular}{lr}
  \hline
Frequency of Posting & Mean Satisfaction \\ 
  \hline
Never & 2.42 \\ 
  Often & 2.02 \\ 
  Rarely & 2.48 \\ 
  Sometimes & 2.42 \\ 
   \hline
\end{tabular}
\end{table}
\begin{table}[ht]
\centering
\caption{Respondents Believing Social Media is Bad
             and Mean Satisfaction by How Often They Post Issues} 
\begin{tabular}{lr}
  \hline
Frequency of Posting & Mean Satisfaction \\ 
  \hline
Never & 2.12 \\ 
  Often & 1.72 \\ 
  Rarely & 2.17 \\ 
  Sometimes & 2.01 \\ 
   \hline
\end{tabular}
\end{table}

\hypertarget{statistical-analysis}{%
\section{Statistical Analysis}\label{statistical-analysis}}

\hypertarget{linear-regression-analysis}{%
\subsection{Linear Regression
Analysis}\label{linear-regression-analysis}}

We will run a linear regression analysis of believing whether social
media is good or bad for democracy on satisfaction with U.S. democracy.

For our first model, we do not include race, ideology, or other
variables in the regression. The p-value for this model is well below
the 5\% threshold, so the null hypothesis can be rejected and the belief
of social media's impact on democracy does have an effect on
satisfaction with democracy. The \(beta_1\) coefficient in the model
(0.3245) suggests a positive correlation between our two variables,
meaning those who believe social media has more of a good effect on
democracy have slightly higher satisfaction ratings of U.S. democracy.

\begin{table}[ht]
\centering
\caption{Model 1} 
\begin{tabular}{rrrrr}
  \hline
 & Estimate & Std. Error & t value & Pr($>$$|$t$|$) \\ 
  \hline
(Intercept) & 2.0918 & 0.0195 & 107.02 & 0.0000 \\ 
  social\_demo & 0.3245 & 0.0316 & 10.26 & 0.0000 \\ 
   \hline
\end{tabular}
\end{table}

We then add in the other factors, beginning with race, then followed by
age category, ideology, and how often the respondent posts social or
political issues on social media.

Adding race:

\begin{table}[ht]
\centering
\caption{Model 2} 
\begin{tabular}{rrrrr}
  \hline
 & Estimate & Std. Error & t value & Pr($>$$|$t$|$) \\ 
  \hline
(Intercept) & 2.0469 & 0.0822 & 24.90 & 0.0000 \\ 
  factor(social\_demo)1 & 0.2933 & 0.0324 & 9.05 & 0.0000 \\ 
  factor(race)2 & 0.1841 & 0.1047 & 1.76 & 0.0788 \\ 
  factor(race)3 & 0.2446 & 0.1089 & 2.25 & 0.0248 \\ 
  factor(race)4 & 0.1649 & 0.0943 & 1.75 & 0.0805 \\ 
  factor(race)5 & 0.0288 & 0.0830 & 0.35 & 0.7288 \\ 
   \hline
\end{tabular}
\end{table}

Race, Ideology, and Age Category:

\begin{table}[ht]
\centering
\caption{Model 3} 
\begin{tabular}{rrrrr}
  \hline
 & Estimate & Std. Error & t value & Pr($>$$|$t$|$) \\ 
  \hline
(Intercept) & 1.9783 & 0.0898 & 22.03 & 0.0000 \\ 
  factor(social\_demo)1 & 0.2980 & 0.0326 & 9.14 & 0.0000 \\ 
  factor(race)2 & 0.1809 & 0.1049 & 1.73 & 0.0846 \\ 
  factor(race)3 & 0.2384 & 0.1090 & 2.19 & 0.0289 \\ 
  factor(race)4 & 0.1548 & 0.0947 & 1.64 & 0.1022 \\ 
  factor(race)5 & 0.0224 & 0.0834 & 0.27 & 0.7883 \\ 
  factor(agecat)2 & 0.0904 & 0.0489 & 1.85 & 0.0646 \\ 
  factor(agecat)3 & 0.0918 & 0.0516 & 1.78 & 0.0757 \\ 
  factor(agecat)4 & 0.0638 & 0.0550 & 1.16 & 0.2465 \\ 
   \hline
\end{tabular}
\end{table}

Race, Ideology, Age Category, and Posting:

\begin{table}[ht]
\centering
\caption{Model 4} 
\begin{tabular}{rrrrr}
  \hline
 & Estimate & Std. Error & t value & Pr($>$$|$t$|$) \\ 
  \hline
(Intercept) & 1.7210 & 0.1072 & 16.06 & 0.0000 \\ 
  factor(social\_demo)1 & 0.2803 & 0.0325 & 8.64 & 0.0000 \\ 
  factor(race)2 & 0.1123 & 0.1026 & 1.09 & 0.2737 \\ 
  factor(race)3 & 0.1781 & 0.1067 & 1.67 & 0.0954 \\ 
  factor(race)4 & 0.0733 & 0.0932 & 0.79 & 0.4315 \\ 
  factor(race)5 & 0.0059 & 0.0818 & 0.07 & 0.9425 \\ 
  factor(agecat)2 & 0.1216 & 0.0482 & 2.52 & 0.0117 \\ 
  factor(agecat)3 & 0.1587 & 0.0513 & 3.09 & 0.0020 \\ 
  factor(agecat)4 & 0.1681 & 0.0550 & 3.05 & 0.0023 \\ 
  factor(ideology)2 & -0.0160 & 0.0387 & -0.41 & 0.6798 \\ 
  factor(ideology)3 & -0.3659 & 0.0407 & -8.99 & 0.0000 \\ 
  factor(posting)2 & 0.3215 & 0.0695 & 4.63 & 0.0000 \\ 
  factor(posting)3 & 0.4284 & 0.0654 & 6.55 & 0.0000 \\ 
  factor(posting)4 & 0.3966 & 0.0659 & 6.02 & 0.0000 \\ 
   \hline
\end{tabular}
\end{table}

\hypertarget{results-discussion}{%
\subsection{Results Discussion}\label{results-discussion}}

\hypertarget{conclusion}{%
\section{Conclusion}\label{conclusion}}

\hypertarget{citations-and-relevant-literature}{%
\section{Citations and Relevant
Literature}\label{citations-and-relevant-literature}}

Dixon, Stacy Jo. Social Media Usage in the United States - Statistics \&
Facts, Statista, 18 Dec.~2023,
\url{https://www.statista.com/topics/3196/social-media-usage-in-the-united-states/\#topicOverview}.

Wike, Richard, et al.~``Social Media Seen as Mostly Good for Democracy
across Many Nations, but U.S. Is a Major Outlier.'' Pew Research
Center's Global Attitudes Project, Pew Research Center, 6 Dec.~2022,
\url{https://www.pewresearch.org/global/2022/12/06/social-media-seen-as-mostly-good-for-democracy-across-many-nations-but-u-s-is-a-major-outlier/}.

\hypertarget{dataset}{%
\subsection{Dataset}\label{dataset}}

American Trends Panel Wave 105
\url{https://www.pewresearch.org/global/dataset/american-trends-panel-wave-105/}

\end{document}
